\sheettitle
  {Exercise sheet}
  {Complete every exercise before opening the independent solutions}

\section*{Exercise 0 -- The $L^p$--RN retention gate}

Let $q>1$, let $Y\in L^q(\mathbb P)$, and let $\mathcal G\subseteq\mathcal F$ be a sub-$\sigma$-field. Define
\[
\nu_Y^{\mathcal G}(A)=\mathbb E[Y\ind_A],
\qquad A\in\mathcal G.
\]
\begin{enumerate}
  \item State why $Y\in L^1(\mathbb P)$ and give the exact norm inequality.
  \item Identify the numerator measure, reference measure, and measurable space. State why the numerator is finite signed and absolutely continuous with respect to the reference measure; do not re-prove countable additivity.
  \item State the signed Radon--Nikodym conclusion, including the space containing the representative and its uniqueness relation.
  \item Let $X$ be a random element with law $\mu_X=X_\#\mathbb P$. Explain why LOTUS requires no assumption $\mu_X\ll\lambda$, whereas a density representation $f_X=d\mu_X/d\lambda$ does.
\end{enumerate}
Limit: 15 minutes. This is a recognition gate; the full Hahn-first construction is not requested.

\section*{Exercise 1 -- Expectation through the law}

Let $X:(\Omega,\mathcal F,\mathbb P)\to(S,\mathcal S)$ be measurable and let $\mu_X=X_\#\mathbb P$.
\begin{enumerate}
  \item Prove for indicators, then for nonnegative simple functions, that
  \[
  \mathbb E[\varphi(X)]
  =
  \int_S\varphi\,d\mu_X.
  \]
  \item Extend the formula to nonnegative measurable $\varphi$ and state the convergence theorem used.
  \item Extend it to signed $\varphi$ whenever the integral is defined.
  \item Give one example of a nonnegative $X$ with $\mathbb E[X]=+\infty$, and one signed $X$ for which $\mathbb E[X]$ is undefined because both parts have infinite expectation.
  \item Explain why discrete sums and continuous density integrals are representations of this identity rather than definitions of different expectations.
\end{enumerate}

\section*{Exercise 2 -- Subgraphs, tail integration, and insurance payments}

Let $(E,\mathcal A,\mu)$ be a $\sigma$-finite measure space and let $f:E\to[0,+\infty]$ be measurable.
\begin{enumerate}
  \item Prove that
  \[
  H_f=\{(x,t)\in E\times(0,+\infty):t<f(x)\}
  \]
  is $\mathcal A\otimes\mathcal B((0,+\infty))$-measurable.
  \item Apply Tonelli to $\ind_{H_f}$ and prove
  \[
  \int_Ef\,d\mu
  =
  \int_0^\infty\mu(f>t)\,dt.
  \]
\end{enumerate}
Now let $X\ge0$ be a random variable and let $d,u>0$.
\begin{enumerate}[resume]
  \item Prove
  \[
  \mathbb E[(X-d)_+]
  =
  \int_d^\infty\mathbb P(X>t)\,dt,
  \qquad
  \mathbb E[X\wedge u]
  =
  \int_0^u\mathbb P(X>t)\,dt.
  \]
  \item If $X\sim\operatorname{Exp}(\lambda)$, compute both expectations.
  \item Determine $\mathbb E[\min\{(X-d)_+,u\}]$ in tail form and then for an exponential loss.
  \item State one convexity regime in which Jensen compares $\mathbb E[g(X)]$ with $g(\mathbb E[X])$, and explain why the comparison does not compute the expected payment.
\end{enumerate}

\section*{Exercise 3 -- Conditioning on a finite partition}

Let $(B_i)_{1\le i\le r}$ be a measurable partition and let $Y\in L^1(\mathbb P)$. Assume first that $\mathbb P(B_i)>0$ for every $i$.
\begin{enumerate}
  \item Determine the general form of a $\sigma(B_1,\ldots,B_r)$-measurable finite real-valued random variable.
  \item Use the conditional-expectation integral identities to prove
  \[
  \mathbb E[Y\mid\sigma(B_1,\ldots,B_r)]
  =
  \sum_i
  \frac{\mathbb E[Y\ind_{B_i}]}{\mathbb P(B_i)}\ind_{B_i}.
  \]
  \item Deduce the law of total expectation.
  \item Treat a partition containing a null atom and identify exactly what remains undetermined.
  \item For an event $A$, apply the formula to $Y=\ind_A$ and recover the law of total probability.
\end{enumerate}

\section*{Exercise 4 -- Bayes and representation bookkeeping}

Let $(B_i)_{1\le i\le r}$ be a partition with positive probabilities and let $A$ satisfy $\mathbb P(A)>0$.
\begin{enumerate}
  \item Derive Bayes' formula from the joint masses $\mathbb P(A\cap B_i)$.
  \item Define $\mathbb P_A(C)=\mathbb P(C\cap A)/\mathbb P(A)$ and prove
  \[
  \frac{d\mathbb P_A}{d\mathbb P}
  =
  \frac{\ind_A}{\mathbb P(A)}.
  \]
  \item For $Y\in L^1(\mathbb P)$ and $\mathcal G=\sigma(B_1,\ldots,B_r)$, fill in the numerator measure, reference measure, domain, and almost-everywhere basis for
  \[
  \frac{d\mathbb P_A}{d\mathbb P},
  \qquad
  \mathbb E[Y\mid\mathcal G].
  \]
  \item Explain why $\mathbb P(C\mid X=x)$ cannot in general be defined by an event ratio when $X$ has a continuous law.
\end{enumerate}

\section*{Exercise 5 -- Original deployment analogues}

The following are original problems. The labels record only the structural mechanism of the corresponding official-question anchors; no official problem statement is reproduced.
\begin{enumerate}
  \item \textbf{Capped-payment mechanism (Q50 anchor).}
  A nonnegative loss $X$ satisfies
  \[
  \mathbb P(X>t)=(1+t)^{-2},
  \qquad t\ge0.
  \]
  Identify the transformed payment and compute $\mathbb E[X\wedge1]$.
  \item \textbf{Nonlinear Poisson moment (Q386 anchor).}
  Let $N\sim\operatorname{Poisson}(2)$ and let the payment be
  \[
  g(N)=\frac{N(N+1)}2.
  \]
  Compute $\mathbb E[g(N)]$ and state why $g(\mathbb E[N])$ is not the required object.
  \item \textbf{Finite-mixture inversion (Q370 anchor).}
  A policy belongs to class $B_1$ with probability $0.3$ and to class $B_2$ with probability $0.7$. An observed event $A$ satisfies
  \[
  \mathbb P(A\mid B_1)=0.8,
  \qquad
  \mathbb P(A\mid B_2)=0.2.
  \]
  Form the two joint masses and compute $\mathbb P(B_1\mid A)$.
\end{enumerate}

\section*{Closed-book reconstruction}

Without notes, complete the two chains
\[
\boxed{X}
\longmapsto
\boxed{\mu_X}
\longmapsto
\boxed{\int\varphi\,d\mu_X}
\]
and
\[
\boxed{\text{finite partition}}
\longmapsto
\boxed{\text{atomwise form}}
\longmapsto
\boxed{\text{integral identities}}
\longmapsto
\boxed{\text{Bayes / total expectation}}.
\]
For every density encountered, name the numerator, reference measure, domain, and almost-everywhere basis. For $\mu_X$, explain why those density fields are not applicable until a reference measure is chosen.

\section*{Deferred proof debts}

The full Hahn-first construction of Radon--Nikodym, $\sigma$-finite gluing, and general conditional kernels or conditional densities are not completion requirements for this packet.
